\section*{Chapitre \ref{ch:b3:environmental-toxicology} --- Chimie de l'environnement et toxicologie}

\begin{solution}{exo:b3:environmental-toxicology:1}
Une molécule diatomique homonucléaire a une seule vibration, qui laisse son
moment dipolaire nul, et l'argon n'a aucune vibration~: aucun ne peut
absorber le rayonnement infrarouge. $\ce{CO2}$ absorbe par
son mode de déformation (et son élongation antisymétrique), qui créent un
dipôle oscillant~; la déformation se situe au milieu de l'émission
terrestre.
\end{solution}

\begin{solution}{exo:b3:environmental-toxicology:2}
$10^{0.1} = 1.26$~: $[\ce{H+}]$ augmente de 26\,\%.
\end{solution}

\begin{solution}{exo:b3:environmental-toxicology:3}
$(2.0 + 0.50)\ \unit{mmol/L} \times \qty{100.1}{mg/mmol} = \qty{250}{mg/L}$ de $\ce{CaCO3}$.
\end{solution}

\begin{solution}{exo:b3:environmental-toxicology:4}
N~: $0.60/14.0 = \qty{0.043}{mmol/L}$~; P~: $0.010/31.0 = \qty{3.2e-4}{mmol/L}$~; N\,:\,P $= 133$, bien au-dessus de 16~: le phosphore s'épuise
le premier et limite la croissance.
\end{solution}

\begin{solution}{exo:b3:environmental-toxicology:5}
$L_0 = \mathrm{BOD_5}/(1 - \eu^{-0.23 \times 5}) = 170/0.683 = \qty{250}{mg/L}$.
\end{solution}

\begin{solution}{exo:b3:environmental-toxicology:6}
$\qty{4.40e-3}{L} \times \qty{0.0200}{mol/L} = \qty{8.80e-5}{mol}$ dans \qty{0.1000}{L}~: \omterm{def:b3:environmental-toxicology:alkalinity}{alcalinité} $\qty{8.80e-4}{mol/L}$~; sous forme de $\ce{HCO3-}$ (\qty{61.0}{g/mol}),
\qty{53.7}{mg/L}.
\end{solution}

\begin{solution}{exo:b3:environmental-toxicology:7}
$K_d = 0.020 \times 300 = \qty{6.0}{L/kg}$. Sorbé/total $= K_dm_s/(K_dm_s + V_w) = 9.0/(9.0 + 0.30) = 0.97$~: 97\,\% sorbé.
\end{solution}

\begin{solution}{exo:b3:environmental-toxicology:8}
Benzène~: $\log\mathrm{BCF} = 0.85 \times 2.13 - 0.70 = 1.11$, BCF d'environ 13. PCB~: $0.85 \times 6.67 - 0.70 = 4.97$, BCF d'environ
\num{9e4}~: il se concentre dans les poissons quelque sept mille fois plus que le benzène~; à de telles valeurs de
$K_{ow}$, la régression approche de sa limite de validité.
\end{solution}

\begin{solution}{exo:b3:environmental-toxicology:9}
$\qty{192}{mg/kg} \times \qty{70}{kg} = \qty{13}{g}$, soit environ 150 tasses. Les espèces diffèrent par l'absorption et le métabolisme,
et une $\mathrm{LD_{50}}$ n'est pas un seuil de sécurité~; les effets nocifs commencent bien au-dessous.
\end{solution}

\begin{solution}{exo:b3:environmental-toxicology:10}
PNEC $= \qty{0.50}{mg/L}/1000 = \qty{0.50}{\micro g/L}$~; $\mathrm{RQ} = 0.20/0.50 = 0.40$, inférieur à 1~: pas de préoccupation pour cette exposition.
\end{solution}

\begin{solution}{exo:b3:environmental-toxicology:11}
Plancton/eau~: $\qty{0.015}{mg/kg}/\qty{2.0e-6}{mg/L} = \qty{7.5e3}{L/kg}$ (bioconcentration). Poisson/plancton~: $0.30/0.015 = 20$~;
oiseau/poisson~: $6.0/0.30 = 20$ (\omterm{def:b3:environmental-toxicology:bio}{bioamplification} à chaque étape)~; oiseau/eau~: $\num{3e6}$.
\end{solution}

\begin{solution}{exo:b3:environmental-toxicology:12}
$\qty{1.0e9}{g}/\qty{64.1}{g/mol} = \qty{1.56e7}{mol}$ de $\ce{SO2}$, ce qui donne autant de moles de $\ce{H2SO4}$~: $\qty{1.53e9}{g}$, environ \qty{1500}{t}. Cela
libère $\qty{3.12e7}{mol}$ de $\ce{H+}$~; à pH 4.0, $[\ce{H+}] = \qty{1.0e-4}{mol/L}$~: $\qty{3.1e11}{L}$ de pluie.
\end{solution}

\begin{solution}{pb:b3:environmental-toxicology:1}
\textbf{1.} Il est dix mille fois plus soluble dans l'octanol que dans
l'eau~: hydrophobe, il se sorbera sur la matière organique et s'accumulera
dans les graisses.

\textbf{2.} $K_{oc} = 0.41 \times 10^4 = \qty{4100}{L/kg}$~; $K_{sw} = f_{oc}K_{oc} = 0.040 \times 4100 = \qty{164}{L/kg}$.

\textbf{3.} $K_{fw} = 0.048 \times 10^4 = \qty{480}{L/kg}$.

\textbf{4.} Non~: à l'équilibre, l'air contient $10^{-4}$ fois la concentration dans l'eau.

\textbf{5.} La molécule neutre et hydrophobe se dissout dans la matière
organique, comme dans l'octanol~; les surfaces minérales sont polaires et
couvertes d'eau.

\textbf{6.} Le sédiment, dont la capacité, $K_{sw}m_s$, est la plus grande.

\textbf{7.} $V_w = \qty{1.0e9}{L}$~; $K_{aw}V_a = \qty{1.0e8}{L}$~; $K_{sw}m_s = \qty{1.64e10}{L}$~; $K_{fw}m_f = \qty{4.8e6}{L}$.

\textbf{8.} $C_w = \qty{1.0e8}{mg}/\qty{1.75e10}{L} = \qty{5.7e-3}{mg/L}$, \qty{5.7}{\micro g/L}.

\textbf{9.} Eau \qty{5.7}{kg}, air \qty{0.57}{kg}, sédiment \qty{94}{kg}, poissons \qty{0.027}{kg}.

\textbf{10.} Air $\qty{5.7e-7}{mg/L}$~; sédiment \qty{0.94}{mg/kg}~; poissons \qty{2.7}{mg/kg}.

\textbf{11.} L'équilibre entre tous les compartiments à la fois, sans
dégradation, sans entrée ni sortie, avec des compartiments parfaitement
mélangés.

\textbf{12.} Avec la dégradation et des flux en régime stationnaire (niveau
II), le stock est fixé par les vitesses d'apport et de perte~; avec des
vitesses de transfert entre compartiments (niveau III), les concentrations
ne sont plus dans les rapports d'équilibre, et le compartiment qui reçoit
l'apport est enrichi.

\textbf{13.} $k = \ln 2/\qty{60}{d} = \qty{0.0116}{d^{-1}}$.

\textbf{14.} $\eu^{-0.0116 \times 365} = 0.0145$~: il en reste 1.5\,\%.

\textbf{15.} $\qty{5.7}{\micro g/L} \times 0.0145 = \qty{0.083}{\micro g/L}$.

\textbf{16.} La dégradation est souvent plus lente dans un sédiment froid et
anoxique, et le pesticide sorbé est relargué lentement dans l'eau, ce qui
le maintient présent après que l'eau s'est éclaircie.

\textbf{17.} Des liaisons qui résistent à l'hydrolyse, à l'oxydation et à
l'attaque microbienne (les liaisons C--Cl aromatiques, par exemple), une
faible solubilité dans l'eau, et une sorption qui le protège des
organismes dégradants.

\textbf{18.} PNEC $= \qty{50}{\micro g/L}/10 = \qty{5.0}{\micro g/L}$.

\textbf{19.} $\mathrm{RQ} = 5.7/5.0 = 1.1$.

\textbf{20.} Juste au-dessus de 1~: un risque est indiqué, et l'évaluation
doit être affinée (meilleures données de toxicité, concentrations
mesurées) ou l'exposition réduite.

\textbf{21.} \qty{2.7}{mg/kg}, presque trois fois le seuil de \qty{1.0}{mg/kg}~: les poissons ne doivent pas être consommés
tant que la concentration n'a pas diminué.

\textbf{22.} Les concentrations dans l'eau, le sédiment et les poissons au
cours du temps et en plusieurs points, pour vérifier le modèle et suivre
la décroissance.

\textbf{23.} Moins de pesticide appliqué, des épandages éloignés des pluies
et des rives, des bandes enherbées qui retiennent le ruissellement, ou une
solution de remplacement moins persistante.

\textbf{24.} Il divise le niveau sans effet mesuré pour couvrir l'écart
entre quelques espèces de laboratoire et tout un écosystème~; un facteur de
10 au lieu de 1000 change la PNEC, et le verdict, d'un facteur cent.

\textbf{25.} Le \omterm{def:b3:environmental-toxicology:ecotoxicology}{quotient de risque} PEC/PNEC vaut environ 1.1, juste au-dessus de 1.
\end{solution}
